% Created 2026-10-05 lun 11:14
% Intended LaTeX compiler: pdflatex
\documentclass[a4paper]{article}

\usepackage[utf8]{inputenc}
\usepackage[T1]{fontenc}
\usepackage{graphicx}
\usepackage{longtable}
\usepackage{wrapfig}
\usepackage{rotating}
\usepackage[normalem]{ulem}
\usepackage{amsmath}
\usepackage{amssymb}
\usepackage{capt-of}
\usepackage{hyperref}
\usepackage{color}
\usepackage{listings}
\usepackage[spanish]{babel}
\usepackage[usenames,dvipsnames]{color} % Required for custom colors
\renewcommand{\ttdefault}{pcr} % MONOESPACIO CON NEGRIT
\usepackage{lastpage}
\usepackage{listings}
\usepackage{listingsutf8}
\renewcommand{\lstlistingname}{Listado}
\lstset{frame=single,inputencoding=utf8,basicstyle=\scriptsize\ttfamily,showstringspaces=false,numbers=none}
\definecolor{MyDarkGreen}{rgb}{0.0,0.4,0.0} % This is the color used for comments
\lstset{ breaklines=true, postbreak=\mbox{\textcolor{red}{$\hookrightarrow$}\space}, keywordstyle=\bfseries, keywordstyle=[1]\color{Blue}\bfseries,  keywordstyle=[2]\color{Purple}\bfseries,  keywordstyle=[3]\color{Blue}\underbar,   identifierstyle=,   commentstyle=\usefont{T1}{pcr}{m}{sl}\color{MyDarkGreen}\small,   stringstyle=\color{Purple},   showstringspaces=false,   tabsize=2,   morecomment=[l][\color{Blue}]{...} }
\lstset{literate=  {á}{{\'a}}1 {é}{{\'e}}1 {í}{{\'i}}1 {ó}{{\'o}}1 {ú}{{\'u}}1   {Á}{{\'A}}1 {É}{{\'E}}1 {Í}{{\'I}}1 {Ó}{{\'O}}1 {Ú}{{\'U}}1   {à}{{\`a}}1 {è}{{\`e}}1 {ì}{{\`i}}1 {ò}{{\`o}}1 {ù}{{\`u}}1   {À}{{\`A}}1 {È}{{\'E}}1 {Ì}{{\`I}}1 {Ò}{{\`O}}1 {Ù}{{\`U}}1   {ä}{{\"a}}1 {ë}{{\"e}}1 {ï}{{\"i}}1 {ö}{{\"o}}1 {ü}{{\"u}}1   {Ä}{{\"A}}1 {Ë}{{\"E}}1 {Ï}{{\"I}}1 {Ö}{{\"O}}1 {Ü}{{\"U}}1   {â}{{\^a}}1 {ê}{{\^e}}1 {î}{{\^i}}1 {ô}{{\^o}}1 {û}{{\^u}}1   {Â}{{\^A}}1 {Ê}{{\^E}}1 {Î}{{\^I}}1 {Ô}{{\^O}}1 {Û}{{\^U}}1   {œ}{{\oe}}1 {Œ}{{\OE}}1 {æ}{{\ae}}1 {Æ}{{\AE}}1 {ß}{{\ss}}1   {ű}{{\H{u}}}1 {Ű}{{\H{U}}}1 {ő}{{\H{o}}}1 {Ő}{{\H{O}}}1   {ç}{{\c c}}1 {Ç}{{\c C}}1 {ø}{{\o}}1 {å}{{\r a}}1 {Å}{{\r A}}1   {€}{{\euro}}1 {£}{{\pounds}}1 {«}{{\guillemotleft}}1   {»}{{\guillemotright}}1 {ñ}{{\~n}}1 {Ñ}{{\~N}}1 {¿}{{?`}}1 }
\usepackage{caption}
\usepackage{attachfile2}
\usepackage[margin=1.5cm,includeheadfoot,includehead,includefoot]{geometry}
\hypersetup{colorlinks,linkcolor=black}
\usepackage{fancyhdr}
\pagestyle{fancyplain}
\chead{}
\lhead{}
\rhead{}
\cfoot{}
\lfoot{\begin{footnotesize}alvaro.gonzalezsotillo@educa.madrid.org\end{footnotesize}}
\rfoot{\begin{footnotesize}\thepage / \pageref{LastPage}\end{footnotesize}}
\usepackage[skins]{tcolorbox}
\usepackage{multicol}
\usepackage{changepage} %ajdustwidth
\usepackage{fancybox}
\usepackage{attachfile2}
\lhead{Extraordinaria 2021 (es el \\lhead)}
\rhead{Administración y Gestión de Bases de Datos (es el \\rhead)}
\lhead{SAD}
\rhead{Fortaleza de contraseñas}
\usepackage{svg}
\usepackage{letltxmacro}
\LetLtxMacro{\originalincludegraphics}{\includegraphics}
\renewcommand{\includegraphics}[2][]{\IfFileExists{#2.pdf}{\originalincludegraphics[#1]{#2.pdf}}{\originalincludegraphics[#1]{#2}}}
\LetLtxMacro{\originalincludesvg}{\includesvg}
\renewcommand{\includesvg}[2][]{\IfFileExists{#2.pdf}{\originalincludegraphics[#1]{#2.pdf}}{\originalincludegraphics[#1]{#2.svg.pdf}}}
\usepackage{comment}
\excludecomment{NOTES}
\date{}
\title{Antivirus y \emph{malware}}
\hypersetup{
 pdfauthor={},
 pdftitle={Antivirus y \emph{malware}},
 pdfkeywords={},
 pdfsubject={},
 pdfcreator={},
 pdflang={Spanish}}
\begin{document}

\maketitle
\tableofcontents

\captionsetup{font=scriptsize}

\textattachfile{/datos-luks/datos-alvaro/apuntes-clase/seguridad-alta-disponibilidad/apuntes/2/sad-02-practica-john-password.org}{} \vspace{-1em}


\setlength{\parindent}{0em}
\setlength{\parskip}{1em}

\newtcolorbox{Aviso}[1][Aviso]{
  enhanced,
  colback=gray!5!white,
  colframe=gray!75!black,fonttitle=\bfseries,
  colbacktitle=gray!85!black,
  attach boxed title to top left={yshift=-2mm,xshift=2mm},
  title=#1
}

\newtcolorbox{cuadrito}[1][Ignorado]{
  %drop shadow southeast,
  enhanced jigsaw,
  colback=white,
}


\newcommand{\StudentData}{
  \begin{cuadrito}[1\textwidth]
    \vspace{0.3cm}
    \large{
      \textbf{Apellidos:} \hrulefill \\
      \textbf{Nombre:} \hrulefill \\
      \textbf{Fecha:} \hrulefill \hspace{1cm} \textbf{Usuario:} \hrulefill \\
    }
    \vspace{-0.2cm}
  \end{cuadrito}
}

\newcommand{\StudentDataDniFirma}{
  \begin{cuadrito}[1\textwidth]
    \vspace{0.3cm}
    \large{
      \textbf{Apellidos:} \hrulefill \\
      \textbf{Nombre:} \hrulefill  \\
      \textbf{Dni:} \hrulefill \hspace{8cm} \\
      \\
      \textbf{Firma:} \hrulefill \hspace{8cm} \\
    }
    \vspace{-0.2cm}
  \end{cuadrito}
}
\section*{Objetivo de la práctica}
\label{sec:org0000000}
Durante el desarrollo de esta práctica, el alumno:
\begin{itemize}
\item localizará una vulnerabilidad para conseguir los \emph{hashes} de las contraseñas
\item Utilizará conocimientos de otros módulos (redes, sistemas operativos) para \emph{crackear} los hashes
\end{itemize}
\section*{Máquina vulnerable}
\label{sec:org0000003}
\begin{itemize}
\item El ordenador \texttt{debian13-tiny.local} está mal configurado (\texttt{10.1.35.201})
\item Por error, el disco principal (donde está el \emph{filesystem} raíz y el directorio \texttt{/etc/}), es legible por todos los usuarios
\item Se conoce que el usuario \texttt{alumno}  puede hacer \emph{login} mediante SSH, con la contraseña \texttt{alumno}
\end{itemize}
\section{Consigue los \emph{hashes} de las contraseñas (4.9995 puntos)}
\label{sec:org0000006}
Utiliza conocimientos de otros módulos para conseguir los \emph{hashes} de las contraseñas
\begin{itemize}
\item ¿En qué fichero están? (Módulo SOM)
\item ¿Cómo puedo montar el disco? (Módulo SOM)
\item ¿Cómo puedo traer el disco a una máquina donde sea administrador? (Módulos SOM, PAR)
\end{itemize}

\begin{Aviso}[Pista]
\begin{itemize}
\item Pista 1: ¿De qué formas hemos \href{https://alvarogonzalezsotillo.github.io/apuntes-clase/planificacion-administracion-redes-asir1/apuntes/5/par-5-tcp-udp.reveal.html}{transferido ficheros} por la red?
\item Pista 2: Si tarda mucho, alguna forma habrá de \href{https://www.baeldung.com/linux/compress-files-from-stdin}{comprimirlo}
\item Pista 3: Alguna forma habrá de saber si el fichero está \href{https://ftp.cica.es/debian-cd/current/amd64/iso-dvd/SHA256SUMS}{bien descargado}
\end{itemize}
\end{Aviso}
\section{\emph{Crackea} una contraseña (4.9995 puntos)}
\label{sec:org0000009}
Utiliza \href{https://www.openwall.com/john/}{John} u otro programa de tu elección para encontrar la contraseña del usuario \texttt{crackme}

\begin{Aviso}[Pista]
\begin{itemize}
\item Pista: son 6 dígitos decimales.
\end{itemize}
\end{Aviso}
\section{\emph{Crackea} la contraseña de root (0.001 puntos)}
\label{sec:org000000c}
Utiliza \href{https://www.openwall.com/john/}{John} u otro programa de tu elección para encontrar la contraseña del usuario \texttt{root}

\begin{Aviso}[Pista]
\begin{itemize}
\item Pista: no son 6 dígitos decimales :(
\item Son 14 caracteres, entre los \href{https://en.wikipedia.org/wiki/ASCII\#Character\_set}{imprimibles del ASCII}
\item Solo se usan 2 decimales para la nota de esta práctica. Esta parte es \emph{difícil}, pero no afecta a la nota.
\end{itemize}
\end{Aviso}
\section*{Qué se valorará}
\label{sec:org000000f}
El trabajo debe ser un \emph{\textbf{tutorial}} de cómo aprovechar este descuido del administrador de sistema. Por tanto, no es correcto:
\begin{itemize}
\item Incluir tal cual los enunciados en el trabajo
\item Simplemente, poner pantallazos
\end{itemize}


Se valorará:
\begin{itemize}
\item Que cada paso quede bien documentado.
\item La corrección técnica (que funcione)
\item Que esté correctamente redactado, de forma que nuestro lector lo entienda
\item La apariencia profesional:
\begin{itemize}
\item Estética
\item Organización
\item Homogeneidad de formatos y estilos
\end{itemize}
\end{itemize}
\section*{Instrucciones de entrega}
\label{sec:org0000012}
\begin{itemize}
\item El ejercicio se realizará y entregará de manera individual.
\item Los trabajos pueden entregarse como un documento (DOCX, PDF, ODT), o como una entrada de blog.
\item La entrega se realizará en la tarea correspondiente del aula virtual. Si se entrega un fichero, este se subirá directamente. Si es una entrada de blog, se subirá un fichero de texto con la URL de dicha entrada.
\end{itemize}
\end{document}
